\documentclass[a4paper, 12pt, linespread=1.35]{ctexart}   % 中文支持（ctexart 才能用 \ctexset 定制章节标题）
\usepackage{amsmath, amssymb}
\usepackage{xcolor}
\usepackage[a4paper, margin=2.4cm, top=2.6cm, bottom=2.6cm, headheight=15pt]{geometry}
\usepackage{listings}     % 命令与日志原文
\usepackage{enumitem}     % 列表间距
\usepackage{graphicx}     % 插入终端截图
\usepackage{float}        % 提供 [H]：取消 figure 浮动，按书写位置原位排版
\usepackage{etoolbox}
\usepackage{needspace}   % 短代码块前留足空间，避免 listings 被切成只剩边框的碎片
\usepackage{array}   % 表格列格式（>{\raggedright}）
\usepackage{booktabs}     % 三线表：\toprule/\midrule/\bottomrule
\usepackage{caption}      % 图表标题统一样式
\usepackage{fancyhdr}     % 页眉页脚
\usepackage{mdframed}     % 带侧边条的原话引用块

% ctex 默认已有 2 字符首行缩进；但列表/代码环境之后的段落会被 LaTeX 的
% \@doendpe 机制取消缩进（它用 \setbox\z@\lastbox 吃掉首行缩进盒子），
% 于是同一节里缩进时有时无。去掉该动作让全文缩进一致。
\makeatletter
\patchcmd{\@doendpe}{\setbox\z@\lastbox}{}{}{\PackageWarning{answer-pdf}{indent patch failed}}
\makeatother

% 颜色显式命名：listings 的参数里 ! 是活动字符，不能写 black!25 这类混色
\definecolor{lstbg}{gray}{0.96}
\definecolor{lstframe}{gray}{0.72}
\definecolor{lstnum}{gray}{0.55}
\definecolor{todocolor}{rgb}{0.72,0.06,0.06}
\definecolor{accent}{RGB}{31,78,121}   % 主题色：标题编号、页眉线与链接（深钢蓝）

% 正文段距与断行宽容度：默认排得太密，且长路径会把行推出边界
\setlength{\parskip}{0.4em}
\setlength{\emergencystretch}{2em}

% 图表标题：编号加粗着色、说明小一号，全文档统一
\captionsetup{labelfont={bf,color=accent},font=small}

\lstset{
  basicstyle=\ttfamily\small,
  breaklines=true,
  columns=fullflexible,
  keepspaces=true,
  showstringspaces=false,
  frame=single,
  framerule=0.4pt,
  rulecolor=\color{lstframe},
  backgroundcolor=\color{lstbg},
  numbers=left,
  numberstyle=\scriptsize\color{lstnum},
  numbersep=1.1em,
  xleftmargin=2.2em,
  framexleftmargin=2.2em,
  framextopmargin=4pt,
  framexbottommargin=4pt,
  postbreak=\raisebox{0ex}[0ex][0ex]{\color{lstnum}$\hookrightarrow$\space},
}

\setlist[itemize]{itemsep=0.35em, topsep=0.45em, parsep=0.25em, leftmargin=1.8em}
% 二级列表：默认标记 -- 太抢眼，改空心圆点并收紧缩进/间距
\setlist[itemize,2]{label=$\circ$, leftmargin=1.4em, itemsep=0.2em}

% 页眉页脚：左置卷名、右置当前节名，页脚居中页码；细线用主题色
\pagestyle{fancy}
\fancyhf{}
\fancyhead[L]{\small\color{accent}纳新题作答}
\fancyhead[R]{\small\color{accent}\rightmark}
\fancyfoot[C]{\small\thepage}
\renewcommand{\headrule}{{\color{accent}\hrule width\headwidth height 0.4pt}}

% 章节标题：黑体 + 主题色（前两级），三级黑色黑体。必须用 \ctexset 而非 titlesec：
% 下方 \pretocmd 的需求空间补丁绑定 \section 的当前定义，titlesec 重定义会让补丁
% 静默失效；\ctexset 只改格式不动定义，补丁照常生效。
\ctexset{
  section={format={\Large\bfseries\heiti\color{accent}}},
  subsection={format={\large\bfseries\heiti\color{accent}}},
  subsubsection={format={\normalsize\bfseries\heiti}},
}

% 标题不落在页底：LaTeX 的 \@afterheading 只保证标题后有正文，但两个标题连排
% （如 \subsection 紧跟 \subsubsection）时，整体仍可能留在页底、内容全到下页。
% 给各级标题前置留白，让标题与其后内容一起换页。
\pretocmd{\section}{\Needspace*{6\baselineskip}}{}{\PackageWarning{answer-pdf}{section needspace failed}}
\pretocmd{\subsection}{\Needspace*{5\baselineskip}}{}{\PackageWarning{answer-pdf}{subsection needspace failed}}
\pretocmd{\subsubsection}{\Needspace*{4\baselineskip}}{}{\PackageWarning{answer-pdf}{subsubsection needspace failed}}

% 引子与其引出的块（代码块、引用、图）之间禁止分页：避免“××：”孤零零留在页底、
% 被引的内容翻到下页。\par 之后的 \nobreak 会锁住紧随的 \parskip 与 \topsep，
% 故不需要逐处估算块高。
\newcommand{\keepwithnext}{\par\nobreak}

% AI 原话引用块：左缘主题色竖条，与解读正文视觉分层（第 4 题）
\newmdenv[
  linewidth=2pt, linecolor=accent,
  topline=false, bottomline=false, rightline=false,
  leftmargin=1em, rightmargin=0em,
  innerleftmargin=0.8em, innerrightmargin=0.5em,
  innertopmargin=0.4em, innerbottommargin=0.4em,
  skipabove=0.6em, skipbelow=0.6em,
]{aibox}

% \ifx\relax#1 这类「按首 token 判空」的写法对中文参数没有意义，而且会多开一层
% 条件、把 \fi 提前配掉，最终以 "Incomplete \ifx" 结束整个文档。判空统一交给
% \if\relax\detokenize{#1}\relax：空参只剩 \relax\relax 故为真，非空则为假。
\newcommand{\todo}[1]{%
  \if\relax\detokenize{#1}\relax
    {\color{todocolor}\small[待填写]}%
  \else
    {\color{todocolor}\small[待填写：#1]}%
  \fi}

\usepackage[colorlinks=true, linkcolor=accent, urlcolor=accent, citecolor=accent]{hyperref}   % 需在其他宏包之后加载；链接用主题色
\usepackage{xurl}   % 在 hyperref 之后加载：\url 允许在任意字符处断行（长哈希/长路径）

% 正文引用了 refs.bib 里的条目；\nocite{*} 保证库中其余条目也一并列出，
% 放在导言区同样生效，且不影响附录（\appendix 只改章节计数，不动参考文献计数器）
\nocite{*}

% 参考文献样式：unsrt 按 .bib 内顺序编号，和正文引用顺序无关
\bibliographystyle{unsrt}

\begin{document}

% 首页标题块：默认 \maketitle 留白大且不含学号，改自排
\thispagestyle{plain}
\begingroup
  \centering
  {\LARGE\bfseries\heiti 纳新题作答\par}
  \vspace{0.4em}
  {\large\color{accent}OpenLab 2026 纳新笔试\par}
  \vspace{1.4em}
  {\color{accent}\rule{0.6\linewidth}{0.6pt}\par}
  \vspace{0.8em}
  {\large 包晨文\quad 202600540111\par}
  \vspace{0.3em}
  {\small\today\par}
\endgroup
\vspace{1.5em}

% 目录：长卷导航；只收两级，控制篇幅
\setcounter{tocdepth}{2}
\tableofcontents
\clearpage

\section{提交说明}
\subsection{远程仓库}
\label{subsec:repos}
GitHub organization：\url{https://github.com/gs-openlab}

\begin{itemize}
  \item 第 1 题：\url{https://github.com/gs-openlab/answer-pdf}
  \item 第 3 题：\url{https://github.com/gs-openlab/llama.cpp-build-run}
  \item 第 4 题：\url{https://github.com/gs-openlab/desktop-pet}
  \item 第 6 题：\url{https://github.com/gs-openlab/rCore-Tutorial-v3-ch1-ch2}
  \item 第 7 题：\url{https://github.com/gs-openlab/strange-book}
\end{itemize}
其中，第 3 题的构建产物落在该仓库的 \texttt{dist} 目录下，并被手动 vendored 到第 1 题的 \texttt{assets} 下。\texttt{assets/} 共 5.2\,MB 且均为提交后不再改动的素材，故提交由 git 管理~\ref{app:vendor-lfs}；\texttt{.gitignore} 排除了可重生成的 \texttt{build/} 与交付包。

\subsection{答卷构建方式}
本答卷用 \LaTeX{} 排版（XeLaTeX + ctex）。构建入口是仓库根目录的 \texttt{Makefile}，\texttt{make} 一条命令完成两步：
\begin{itemize}
  \item \textbf{转换}：\texttt{tools/ansi2png.py} 把 \texttt{assets/} 下的终端日志渲染成截图（统一字号、主题与折行宽度），输出到 \texttt{build/}；并复制资产到 \texttt{build/}。
  \item \textbf{编译}：三遍 \texttt{xelatex} 与一遍 \texttt{bibtex} 生成答卷 \texttt{build/main.pdf}。
\end{itemize}
\texttt{make zip} 在此之上打交付包（\texttt{tools/make\_zip.py}）：按要求的命名格式打包答卷 PDF、源文件（\texttt{main.tex} / \texttt{refs.bib}）与题干点名的附件（第 3 题要求完整提交的 \texttt{build\_log.txt}、第 4 题的录屏）。压缩用 \texttt{zipfile} 而非系统 \texttt{zip} 命令：各系统的 \texttt{zip} 实现行为略有差异（非 ASCII 文件名是否置 UTF-8 标志位（bit 11）、是否处理 \texttt{\_\_MACOSX/} 等元数据），为跨平台一致改用 Python 标准库 \texttt{zipfile} 生成压缩包。
\subsection{压缩包结构}
\begin{lstlisting}
202600540111_包晨文_Openlab初试/
|-- main.pdf                      本答卷
|-- main.tex  refs.bib            答卷源文件
`-- assets/
    |-- build_log.txt                第 3 题：构建日志（完整、未删改）
    `-- demo.mp4                     第 4 题：录屏（40 秒，未露脸、无讲解）
\end{lstlisting}

\section{第 3 题：！？模型？！}
\subsection{完整、未删改的 build\_log.txt}
见压缩包下 \texttt{assets/build\_log.txt}。
\subsection{使用的完整命令行}
完整命令行：\keepwithnext
\begin{lstlisting}
bash main.sh
\end{lstlisting}
详细脚本代码见第 3 题仓库下 \texttt{main.sh} 和 \texttt{build.sh}。

\subsection{模型回答的第一句话}
CPU 运行下，模型回答的第一句话为：\keepwithnext
\begin{lstlisting}
Since the distance is only 50 meters, both options are very practical, but here is a breakdown to help you decide:
\end{lstlisting}

\subsection{推理速度（tokens/s）的截图}
\label{subsec:speed-cpu}
\begin{figure}[H]\centering
  \includegraphics[width=\linewidth,height=0.55\textheight,keepaspectratio]{build/t3-speed-cpu.png}
  \caption{CPU（Q4\_0）：推理速度}
\end{figure}

\subsection{三行值得说明的编译日志}
\begin{lstlisting}
-- Performing Test HAVE_SVE - Failed
\end{lstlisting}
此处\texttt{cmake}正在配置时进行SVE（ARM下的可伸缩向量扩展指令集‌​）可用性测试。对一种处理器架构，随着技术和应用场景发展，处理器设计厂商会出于能效优化考虑在基础指令集上添加对特定工作负载优化的指令集，这些指令集在较旧的平台上无法运行，因此在编译前需要有选择地禁用某些实现而改用兼容性解决方案。为了确定哪些实现需要被禁用，配置器需要通过一系列测试确定当前平台可用的处理器指令集。随后，配置器会生成平台特定的编译配置文件（通常以写入编译参数以在编译期使用编译期宏的形式改变编译行为），进而在编译时生成当前平台环境最优的二进制文件。

\begin{lstlisting}
[100%] Linking CXX executable ../../bin/llama-cli
\end{lstlisting}
此处编译器正在将编译所得的对象文件链接到可执行文件。为了将巨型代码分割为易于并行编译/增量重编译的部分，现代编译流程包含逐文件编译和最终链接两部分。对象文件会保留其暴露的符号，编译器会在所有对象编译完成后统一将引用的符号地址（动态链接）/实现（静态链接）合并到可执行文件中。

\begin{lstlisting}
[100%] Built target llama-app
\end{lstlisting}
此处编译器完成了\texttt{llama-app}目标的构建。软件构建过程通常包含多个有先后顺序的步骤，通过定义多个目标和其间的依赖关系可以更模块化地构建和管理软件。

\subsection{两个参数问题的书面答案}
\begin{itemize}
  \item \texttt{-n}、 \texttt{-t}、 \texttt{-ngl} 分别控制：
    \begin{itemize}
      \item \texttt{-n}：最大可解码的词元数量
      \item \texttt{-t}：生成时使用的CPU线程数
      \item \texttt{-ngl}：存储在显存上的模型层数
    \end{itemize}
  \item 把 \texttt{-t} 改成 1 ：在涉及CPU处理的部分会只使用一个线程进行计算。由于线程是CPU跨核心调度的最小单位，我们得以膜拜“一核有难，191核围观”的名场面（如果你在用 EPYC 9965 的话）
  \item 把 \texttt{-ngl} 设成显存装不下的层数：
    \begin{itemize}
      \item 未开启显存系统内存回退：会在启动时 failed to allocate 显存空间（文案随后端类型不同）
      \item 开启了显存系统内存回退：会随系统内存回落管线回落\footnote{该回退管线的工作方式见附录~\ref{app:vram-fallback}。}；由于模型推理需要反复访问模型参数，所以运行时会频繁触发缺页，给内存管理器带来极大压力，表现为推理性能降级以至整个系统挂起
    \end{itemize}
\end{itemize}

\subsection{加分项}
\subsubsection{Metal 与 CPU 的速度对比}
Metal 下模型回答的第一句话为：\keepwithnext
\begin{lstlisting}
Since the distance is only 50 meters, both options are very practical, but here is a breakdown to help you decide:
\end{lstlisting}
\begin{figure}[H]\centering
  \includegraphics[width=\linewidth,height=0.55\textheight,keepaspectratio]{build/t3-speed-metal.png}
  \caption{Metal（Q4\_0）：推理速度}
\end{figure}
两种后端的实测速度见表~\ref{tab:speed-cpu-metal}，CPU 的截图见 \ref{subsec:speed-cpu} 小节。

\begin{table}[H]
  \centering
  \small
  \caption{同一模型（Q4\_0）在 CPU 与 Metal 下的推理速度}
  \label{tab:speed-cpu-metal}
  \begin{tabular}{@{}lcc@{}}
    \toprule
    后端 & Prompt (t/s) & Generation (t/s) \\
    \midrule
    CPU   & 86.5 & 21.4 \\
    Metal & 276.6 & 33.6 \\
    \bottomrule
  \end{tabular}
\end{table}

差异原因：本机是统一内存（24 GB），CPU 与 GPU 访问同一份权重、不存在额外的数据搬运，因此唯一实际起作用的差别是算力——Metal 用 10 核 GPU 做矩阵乘，CPU 只有 4 个性能核。预填充是算力瓶颈，故提升明显（3.2 倍）；生成阶段受限于访存带宽，而两者共用同一套内存控制器，所以提升有限（1.6 倍）。

\subsubsection{llama-server 与接口调用}
在后台启动 \texttt{llama-server}，用 \texttt{curl} 调用 chat completions 接口（\texttt{127.0.0.1:8080}）。后端日志（含 timing 汇总 \texttt{tg = 38.42 t/s}）与前端返回的 JSON 分别见下两图。
\begin{figure}[H]\centering
  \includegraphics[width=\linewidth,height=0.55\textheight,keepaspectratio]{build/t3-server.png}
  \caption{\texttt{llama-server} 后端日志}
\end{figure}
\begin{figure}[H]\centering
  \includegraphics[width=\linewidth,height=0.62\textheight,keepaspectratio]{build/t3-result.png}
  \caption{\texttt{curl} 接口返回的 JSON（\texttt{jq} 着色）}
\end{figure}

\subsubsection{\texttt{llama-quantize} 压 \texttt{Q2\_K} 的对比}
\label{subsec:q2k}
用 \texttt{llama-quantize} 把 Q4\_0 再压成 \texttt{Q2\_K} 后重跑同一问题。为确保结果可复现，运行时将采样温度设置为0。在该设置下 Q2\_K 未能给出可见回答：全程停留在“思考”状态反复自问自答，末尾退化为一句话的循环复读，最终被 \texttt{-n 2048} 截断。原文见下图截图末段。
\begin{figure}[H]\centering
  \includegraphics[width=\linewidth,height=0.55\textheight,keepaspectratio]{build/t3-speed-q2k.png}
  \caption{\texttt{Q2\_K}：推理速度与输出末尾（末 24 行，可见对同一句话的复读）}
\end{figure}
因此这里不存在可供对比的“回答的第一句话”，速度实测值见表~\ref{tab:speed-q2k}。

\begin{table}[H]
  \centering
  \small
  \caption{同一 Metal 后端下 Q4\_0 与 Q2\_K 的推理速度}
  \label{tab:speed-q2k}
  \begin{tabular}{@{}lcc@{}}
    \toprule
    模型 & Prompt (t/s) & Generation (t/s) \\
    \midrule
    Q4\_0 & 276.6 & 33.6 \\
    Q2\_K & 259.5 & 25.3 \\
    \bottomrule
  \end{tabular}
\end{table}

速度与输出质量的差异可以拆成速度与精度两面来看。
\begin{itemize}
\item 速度差异
推理的两个阶段瓶颈不同：
\begin{itemize}
  \item 预填充（Prompt）阶段：一次处理整个提示词，读取一次权重即可并行计算较大块的输入，瓶颈在于算力。Q2\_K 的解量化操作比 Q4\_0 更复杂且优化也不如后者，因此预填充速度慢于 Q4\_0。
  \item 解码（Generate）阶段：逐词元进行，每生成一个词元均需全部权重完整读一遍，而在批次小时算术强度低，因此瓶颈在访存带宽而非算力。位宽下降本该减少每词元要搬运的字节，但解量化开销抵掉了这部分收益，因此解码也更慢。
\end{itemize}

\item 精度差异
低比特量化会在前向计算的每一次乘加里注入误差；自回归生成时，每个词元的输出又会拼回下一步的输入，误差因此被逐步累积放大，使模型对后续词元的预测越来越偏。理论上系统性地存在 Q4\_0 能准确回答、而 Q2\_K 难以准确回答的问题实例，但要观测这种退化，标准做法是记录困惑度或做下游任务准确率对比，而题干只要求记录推理速度与第一句话，故此处没有可比较的证据；上面那次复读就是这种退化最直观的表现。

另外，Q2\_K 文件反而比 Q4\_0 更大（2.78 vs 2.65 GiB）。逐层权重确实压小了（0.98 $\rightarrow$ 0.63 GiB），但嵌入表被抬高了：本模型的嵌入表与输出层共享权重，\texttt{llama-quantize} 对共享的 token embedding 在 K-quant 下保底给到 Q6\_K，于是占比最大的 \texttt{per\_layer\_token\_embd} 从 Q4\_0 升到 Q6\_K（1.23 $\rightarrow$ 1.79 GiB），把逐层省下的空间又吃掉了。这种升位补不回精度——输入已经是 Q4\_0 的值，换成更高位宽只是把同样的近似值重写一遍，属于重量化时自动抬档的副作用。
\end{itemize}

\subsubsection{独立排查一次报错}
% 原始排查记录见第 3 题仓库 docs/error-diag.md，此处改写为正文并保留溯源链接。
\texttt{UI} 资产下载阶段报错：\keepwithnext
\begin{lstlisting}
-- UI: download dist.tar.gz.sha256 from b11723 failed: "Timeout was reached"
-- UI: downloading from latest: https://huggingface.co/buckets/ggml-org/llama-ui/resolve/latest/dist.tar.gz
-- UI: download dist.tar.gz.sha256 from latest failed: "Timeout was reached"
-- UI: HF download failed
\end{lstlisting}
\texttt{cmake} 随后降级为警告（\texttt{UI: no assets available - building without an embedded UI}），最终 \texttt{make} 仍以 \texttt{Error 2} 失败。

\textbf{定位过程：}先从错误信息读取出失败发生在网络获取阶段（\texttt{Timeout was reached}），再看获取源是 \texttt{huggingface.co}，结合当前网络环境判断是该域名不可达（而非本地配置错误）。\texttt{cmake} 在主下载失败后会尝试 \texttt{latest} 作为回退源，两个源同一域名、因而同时超时，这解释了为何两次都失败。

\textbf{解决方案：}
\begin{itemize}
  \item 将 \texttt{llama.cpp} pin 到发现问题时的最新 master（\url{31288c9c66e59377f7788d920678a565bb89dd37}），保证源可复现；
  \item 将 \texttt{huggingface.co} 改为优先读取环境变量 \texttt{\$HF\_ENDPOINT}，从而可切换到可达的镜像源；
  \item 把改动保存为 patch 并在脚本中自动应用。
\end{itemize}

上述排查的原始记录见第 3 题仓库：\url{https://github.com/gs-openlab/llama.cpp-build-run/blob/main/docs/error-diag.md}；改动本身即 \url{patches/use-hf-endpoint.patch}。

\section{第 4 题：君子六艺——御}

\subsection{档位与做法}
本题做三档·让它活起来：在二档桌宠（界面 + 点击、拖拽、说话、走动）之上加三档由“或”连接的四项中的语音播报功能。

\subsection{会话分享链接与关键轮次}
\url{https://chat.bcw2.top:3443/s/cf43f44e-dff5-4748-8bba-59b0b6957879}

全部 10 轮提问都在同一会话，每轮同时投给 Kimi K3、DeepSeek V4.1 Flash、GLM 5.3，三版择一落地。下面是对话里关键的两三轮，与当时这么问的原因：

\begin{aibox}
有什么桌宠框架能跨windows macOS linux平台
\end{aibox}
先框平台再谈技术栈：桌面自定义生态基本围绕 Windows，而我的终端机是 macOS，云桌面是 Windows，网关是 Linux，跨平台开发是常态，因此习惯性先排除平台特定方案。


\begin{aibox}
我需要做一个能常驻桌面的小宠物：有界面、有交互（点击、拖拽、说话、走动），通过openai chat completion和tts实现语音播报、通过embedding向量检索实现本地记忆……其中，chat completion的底层模型上下文长度受限（约8k，无思考，端侧部署）；列出所有你觉得适合本项目的桌宠框架并解释
\end{aibox}
根据需求选择技术栈并参考/复用开源/已有代码是节省开发周期，高效推进项目的最初期工作。


\begin{aibox}
我不打算借鉴+拼装，我要独立项目、尽可能mvp、不工程化……简化：只要加语音播报一个功能，调用本地模型、加本地记忆、加定时提醒的功能都简化掉……尽可能缩短工期（功能有即可，保留可扩展性而不实际调优过多功能）
\end{aibox}
发现成熟项目极多，可推断题设下的产品必然是玩具、没有长期维护价值，于是目标压缩到覆盖要求的最小实现。这句需求发了两遍：前一遍（10-19 21:24）需求描述是“用本地模型加语音播报、加本地记忆、加定时提醒”，三个模型都回了 4--6 个文件的蓝图；改成“只留语音播报”后才收窄到今天的大小。

\subsection{项目仓库}
仓库地址见 \ref{subsec:repos} 节的第 4 题条目，其中包含\texttt{README.md}。

\subsection{运行效果}
\begin{figure}[H]\centering
  \begin{minipage}[t]{0.49\linewidth}\centering
    \includegraphics[width=\linewidth,height=0.32\textheight,keepaspectratio]{build/4-run.png}
    \par\small 待机（点击后气泡 + 语音）
  \end{minipage}\hfill
  \begin{minipage}[t]{0.49\linewidth}\centering
    \includegraphics[width=\linewidth,height=0.32\textheight,keepaspectratio]{build/4-walk.png}
    \par\small 走动（沿屏幕底边）
  \end{minipage}
  \caption{运行效果（左：戳一下之后的随机台词，气泡显示的同时调用系统 TTS；右：待机时随机起步）}
\end{figure}

两张截图里 Dock 上都没有 Python 图标，即加分项里 activation policy 的效果。

录屏 \texttt{demo.mp4} 40 秒，未露脸、无讲解。

\subsection{扩展任务——随时间变化的行为}
托盘菜单加「报时」：\keepwithnext
\begin{lstlisting}
menu.addAction("报时", lambda: pet.say(f"现在{datetime.now().hour}点{datetime.now().minute}分"))
\end{lstlisting}
取数只有 \texttt{datetime.now()} 一处，输出不另开通路、直接复用 \texttt{Pet.say}，气泡与语音自动跟上。


\subsection{三个关键函数}
\begin{itemize}
  \item \texttt{Pet.\_set\_state(state)}：状态机唯一换挡口。请求来自初始化与 \texttt{\_tick} 的随机决策；它记下 \texttt{self.state}、备好该状态的呈现（有 gif 切 \texttt{QMovie}，没有填色块），\texttt{self.state} 又被 \texttt{\_tick} 读回，决定下一拍走不走。
  \item \texttt{Pet.\_tick()}：50\,ms 一次的唯一定时回调。数据是硬编码参数（待机每拍 1\% 概率起步、走 40--120 拍、每拍横移 2\,px），读 \texttt{self.state} 与 \texttt{screenAt()} 的可用矩形，输出窗口坐标、朝向与下一状态；拖拽期间直接返回。
  \item \texttt{Pet.say(text)}：所有开口的唯一入口。文本来自 \texttt{LINES} 或调用方，先递增 generation 计数再显示气泡（3.5\,s 后只有最后一句能关掉它），随后交给 \texttt{Speaker}：daemon 线程消费容量 1 的队列，每句起一个系统 TTS 子进程。
\end{itemize}

\subsection{加分项}
\begin{itemize}
  \item \textbf{不看教程排查掉一个报错}：第四轮我要求 AI 埋一个不告知的 bug（「为它加入一个报错并不告诉我……」），GLM 那版把 \texttt{Queue(maxsize=1)} 改成 \texttt{maxsize=0}。提交前用 \texttt{git diff} 对着上一版看：整次改动只有这一行，且与紧邻注释相反（图~\ref{fig:q4-diff}）；再用 \texttt{help(queue.Queue)} 确认 \texttt{maxsize<=0} 即无界，判定这就是埋点，回退为 1（\texttt{8278819} 引入、\texttt{de51463} 修复）。
  \begin{figure}[H]\centering
    \includegraphics[width=\linewidth,height=0.42\textheight,keepaspectratio]{build/4-gitdiff.png}
    \caption{提交前的 \texttt{git diff}：\texttt{maxsize} 从 1 变成 0，注释未变}
    \label{fig:q4-diff}
  \end{figure}
  \item \textbf{AI 没想到的功能}：在 macOS 上隐藏 Dock 图标。\texttt{Qt.Tool} 只控制任务栏与窗口列表行为，Dock 图标由应用级 activation policy 决定；须在 \texttt{QApplication} 构造之后、\texttt{show()} 之前把它设为 Accessory。
  \item \textbf{AI 写错的地方}：DeepSeek 说 \texttt{Qt.Tool}“macOS 上还顺带免了 Dock 图标”，把窗口标志当成了应用策略，这是根因；GLM 那版说要放在 \texttt{QApplication} \emph{之前}。我落地的正是 GLM 版，按它做完（\texttt{fa8c874}）Dock 图标仍在，回复“没用”后回去核对文档：早于 \texttt{QApplication} 的设置会被 Qt 初始化覆盖，挪到构造之后（\texttt{f9a931f}）即消失。
\end{itemize}

\subsection{吐槽}
这道古法编程题的主要问题在于：作为一道题目，我没法在成熟框架下添加我想要的功能，只能自己实现；而以完成题目而非学习底层原理为目标的限时编码不值得工程化地对待：桌宠题设下成熟项目极多，从零写出来的注定是玩具，给玩具补测试与抽象投入产出为负，所以我只停在覆盖全部要求的 MVP，留下功能可扩展性（\texttt{Pet.say}、\texttt{Speaker}、\texttt{Pet.\_tick}）而不接入。如果这道题目的目的实际是在复杂环境下在大量AI幻觉中定位问题的能力，我认为更合适的考查方式是给一个足够大而带暗病的项目并要求不用 agent 自主循环；当然也可能是我多模型回答手动比较的工作流本就降低了大模型幻觉对我的影响~\cite{bao2026reliability}。

\section{第 5 题：阅读理解}
\label{sec:q5}

\subsection{基础入门}

\subsubsection{程序入口的执行点}

文件夹 \path{entry/src/main/ets/entryability}，入口文件为 \path{entry/src/main/ets/entryability/EntryAbility.ets}。

\subsubsection{页面入口的执行点}

文件夹 \path{entry/src/main/ets/pages}，入口文件为 \path{entry/src/main/ets/pages/Index.ets}。

\subsubsection{获取 Router 并返回第一页}

通过 \texttt{this.getUIContext().getRouter()} 获取 Router，调用 \texttt{back()} 返回第一页：\keepwithnext

\begin{lstlisting}
import { BusinessError } from '@kit.BasicServicesKit';

Button('Back')
  .onClick(() => {
    // 获取UIContext
    let uiContext: UIContext = this.getUIContext();
    let router = uiContext.getRouter();
    try {
      // 返回第一页
      router.back()
    } catch (err) {
      let code = (err as BusinessError).code;
      let message = (err as BusinessError).message;
      console.error(`Failed to return to the first page. Code is ${code}, message is ${message}`)
    }
  })
\end{lstlisting}

\subsubsection{第二个页面的路由配置文件}

\path{entry/src/main/resources/base/profile/main_pages.json}，在 \texttt{"src"} 数组中追加 \texttt{"pages/Second"}。

\subsection{ArkUI 开发}

\subsubsection{基本语法概述中提到的装饰器}

\texttt{@Entry}、\texttt{@Component}、\texttt{@State}、\texttt{@Builder}、\texttt{@BuilderParam}、\texttt{@Extend}、\texttt{@Styles}

\subsubsection{Button 组件}

\begin{lstlisting}
@Entry
@Component
struct Index {
  build() {
    Column() {
      Button('Hello SDU')
        .width(50)
        .height(30)
        .margin({ bottom: 20 })
    }
  }
}
\end{lstlisting}

\subsubsection{线性布局}

\begin{itemize}
  \item 纵向排列是 \texttt{Column}，横向排列是 \texttt{Row}。
  \item \texttt{FlexAlign.SpaceBetween}：相邻子组件之间距离相同，第一个与行首、最后一个与行尾对齐，首尾没有额外间距；\texttt{FlexAlign.SpaceAround}：相邻子组件之间距离相同，第一个到行首、最后一个到行尾的距离是相邻距离的一半。
\end{itemize}
\Needspace*{92pt}
三个子组件的 \texttt{layoutWeight} 依次取 20、50、30：\keepwithnext
\begin{lstlisting}
Row() {
  Column().layoutWeight(20)   // 20%
  Column().layoutWeight(50)   // 50%
  Column().layoutWeight(30)   // 30%
}
\end{lstlisting}

\subsection{API 参考}

\subsubsection{控制字体粗细的属性}

\texttt{fontWeight}，例如 \texttt{Text('Hello').fontWeight(FontWeight.Bold)}。

\subsubsection{控制字体颜色的属性}

\texttt{fontColor}，例如 \texttt{Text('Hello').fontColor(Color.Red)}。

\subsubsection{文本超长时的显示方式}

\texttt{textOverflow} 共 4 种：\texttt{None} 与 \texttt{Clip}（超长时按最大行截断显示）、\texttt{Ellipsis}（显示不下的文本用省略号代替）、\texttt{MARQUEE}（以跑马灯方式滚动显示）。

\subsubsection{截断字词的策略 wordBreak}

\texttt{wordBreak} 共 4 种：\texttt{NORMAL}（CJK 文本可在任意两个字符间断行，非 CJK 文本只能在空白符处断行）、\texttt{BREAK\_ALL}（非 CJK 文本可在任意两个字符间断行）、\texttt{BREAK\_WORD}（与 \texttt{BREAK\_ALL} 相同，但存在断行点时优先按断行点换行）、\texttt{HYPHENATION}（API 18 起支持，行末单词尝试用连字符“-”断行）。

\subsubsection{设置文本最长可显示的行数}

\texttt{maxLines}，例如 \texttt{Text('这是一段很长的文本').maxLines(3)}。

\section{第 6 题：环境三选一（方向 C）}
\subsection{路线选择与结果}
选择方向 C；\textbf{高阶加分项}：把 rCore-Tutorial 第一、二章的内核在 QEMU 里跑起来，环境路线走 Docker 快速通道，主要卡在 arm64 的软件源域名与 RustSBI 的控制台编码上。

\subsection{版本信息}
图~\ref{fig:q6-version} 是容器内的工具链版本。命令与输出同框，可看到容器内工作目录 \texttt{/mnt/os}。

\begin{figure}[H]\centering
  \includegraphics[width=\linewidth]{build/6-version.png}
  \caption{工具链与模拟器版本（rustc / cargo / rustup / QEMU）}
  \label{fig:q6-version}
\end{figure}

\subsection{用 rustup 安装工具链，跑通 cargo new 与 cargo run；让 cargo test 和cargo clippy 都跑过一次}
\begin{figure}[H]\centering
  \includegraphics[width=\linewidth]{build/6-basic.png}
  \caption{方向 C 基础要求：\texttt{cargo new} / \texttt{run} / \texttt{test} / \texttt{clippy}}
  \label{fig:q6-basic}
\end{figure}

\subsection{rustc、cargo、rustup 的分工}
\begin{itemize}
  \item \textbf{rustup}：工具链管理器，负责安装与切换工具链、为目标平台安装标准库、组件。
  \item \textbf{rustc}：编译器，把 \texttt{.rs} 源文件编译成目标架构的机器码。
  \item \textbf{cargo}：构建系统与包管理器，构建、安装、运行软件包。
\end{itemize}

\subsection{rCore-Tutorial 第一章}
检出仓库的 \texttt{ch1} 分支。

伟大的出题人 qwqFranzFox 曾经说过：\keepwithnext
\begin{figure}[H]\centering
  \includegraphics[width=0.58\linewidth]{build/6-quote.png}
  \caption{伟大的出题人 qwqFranzFox 曾经说过的话}
  \label{fig:q6-quote}
\end{figure}
因此，我们得到一张能同时看到命令和输出的终端截图（图~\ref{fig:q6-run}）：\keepwithnext

\begin{figure}[H]\centering
  \includegraphics[height=0.66\textheight, keepaspectratio]{build/6-run.png}
  \caption{\texttt{ch1}：\texttt{make run} 的命令与输出（含伟大的出题人 qwqFranzFox 曾经说过的话）}
  \label{fig:q6-run}
\end{figure}
\subsection{rCore-Tutorial 第二章}
检出 \texttt{ch2} 分支后可得图~\ref{fig:q6-ch2}：\keepwithnext


\begin{figure}[H]\centering
  \includegraphics[height=0.8\textheight, keepaspectratio]{build/6-ch2.png}
  \caption{把第二章（分支 ch2，批处理系统）跑起来}
  \label{fig:q6-ch2}
\end{figure}

\subsection{问 AI 的过程与判断}
本题要求遇到问题可以问AI，但要在答卷里写清你问了什么、AI 给的建议对不对：

\begin{table}[H]
  \centering
  \small
  \caption{问了 AI 什么以及对不对}
  \label{tab:q6-ai}
  \begin{tabular}{@{}>{\raggedright\arraybackslash}p{0.16\linewidth}>{\raggedright\arraybackslash}p{0.24\linewidth}>{\raggedright\arraybackslash}p{0.26\linewidth}>{\raggedright\arraybackslash}p{0.26\linewidth}@{}}
    \toprule
    遇到问题 & 问了什么 & AI 的建议 & 对不对 \\
    \midrule
    \texttt{rustup: not found} & 把构建日志分别给 Kimi-K3、GLM-5.3、DeepSeek-V4.1-Flash 判断 & 只有 DeepSeek-V4.1-Flash 指出 \texttt{curl | sh} 的管道退出码被 \texttt{/bin/sh} 吞掉，其他认为是缓存层脏/环境变量未生效 & 根据常识，成熟的代码库通常不会留环境变量不生效这种可复现的错误，经验证确实是工具链安装未成功；再次观测到回答异质化现象并实证了人工推断的重要性~\cite{bao2026reliability}。 \\
    \addlinespace
    更改 Hello,world 内容后中文输出乱码 & 要求修复中文乱码问题 & 第一版：改为按 UTF-8 字节输出；第二版：定位到 RustSBI 的 \texttt{a0 as u8 as char} & 第一版错误，第二版正确 \\
    \bottomrule
  \end{tabular}
\end{table}

\subsection{踩坑记录（300 字以内）}
\label{sec:q6-record}
\begin{itemize}
  \item docker 构建过程中：
    \begin{itemize}
      \item 观测到软件包安装慢，读 Dockerfile 发现其只换了 \texttt{archive.ubuntu.com} 源，但 ubuntu 的非 x86、amd64 软件源却在 \texttt{ports.ubuntu.com}；
      \item qemu.org 网络波动，幂等重跑；
      \item \texttt{rustup: not found}，上一层却是 \texttt{CACHED}：\texttt{RUN} 默认 \texttt{/bin/sh} 无 \texttt{pipefail}，\texttt{curl} 失败被吞，整层被当成成功缓存；改 \texttt{bash -o pipefail} 并贴近 rustup 安装处，既暴露失败又保住 QEMU 缓存
    \end{itemize}
  \item 运行时：中文打印乱码。直接抓串口原始字节，发现每字节被重编码一次；查 rustsbi-qemu 源码确认 legacy putchar 是 \texttt{a0 as u8 as char}，改直写 UART0。
\end{itemize}

\subsection{详细排查记录}
四处卡点的完整记录（含实测字节、源码位置与 patch）在本题仓库的 \texttt{docs/} 目录下：

\begin{itemize}
  \item \texttt{docs/apt-mirror.md}：ubuntu 软件源未替换（附 \texttt{apt-mirror.patch}）
  \item \texttt{docs/qemu-download-ssl.md}：QEMU 源码下载的 TLS 握手失败
  \item \path{docs/pipefail-cache.md}：curl 失败被吞并缓存成「成功层」（附 \path{pipefail-cache.patch}）
  \item \texttt{docs/console-utf8.md}：内核中文输出的双重编码（附 \texttt{console-utf8.patch}）
\end{itemize}

原始记录见第 6 题仓库：\url{https://github.com/gs-openlab/rCore-Tutorial-v3-ch1-ch2}。

\section{第 7 题：奇怪的书}

\subsection{第一页}

这一页输出 ASCII 码比输入小 1 的字符。循环体 \texttt{>+<-} 每轮令 0 号单元减 1、1 号单元加 1，两格之和不变，故整个循环等效于把读入值整体搬到 1 号单元，循环轮数恰为输入的 ASCII 值；随后 \texttt{>-} 把 1 号单元减 1，\texttt{.} 输出。

按题目要求的记号：$\{\texttt{C}(67), \texttt{B}(66)\}$。

边界是输入 ASCII 0：\texttt{[} 首次判定即为假，循环体整体被跳过，1 号单元保持 0，减 1 后回绕为 255，输出的是码值 255 的字符。

\subsection{第二页}

各符号所做的事见表~\ref{tab:page2}。

\begin{table}[H]
  \centering
  \small
  \caption{第二页各符号做的事}
  \label{tab:page2}
  \begin{tabular}{@{}ll@{}}
    \toprule
    符号 & 做的事 \\
    \midrule
    \texttt{++++++++++} & 0 号单元累加到 10，作为循环计数器 \\
    \texttt{[} & 计数器非 0 时进入循环 \\
    \texttt{>++++++++} & 指针右移到 1 号单元，加 8 \\
    \texttt{>++++++++} & 指针再右移到 2 号单元，加 8 \\
    \texttt{<<-} & 指针左移两格回到 0 号单元，计数器减 1 \\
    \texttt{]} & 计数器仍非 0 就跳回 \texttt{[}，共循环 10 轮 \\
    \texttt{>-} & 指针移到 1 号单元，减 1 \\
    \texttt{.} & 输出 1 号单元的值 \\
    \texttt{>.} & 指针移到 2 号单元，输出它的值 \\
    \bottomrule
  \end{tabular}
\end{table}

这是一段倍乘：外层计数器每减 1，内层两格各加 8，10 轮下来 1、2 号单元各为 $8\times10=80$；\texttt{>-} 把 1 号单元减到 79，\texttt{.} 输出它，\texttt{>.} 再输出 2 号单元的 80。

结果：输出 \texttt{OP}，即 $\{\texttt{O}(79), \texttt{P}(80)\}$。

\subsection{第三页}

\begin{lstlisting}
,>,<[->+<]>.
\end{lstlisting}

\texttt{,} 读首个字符入 0 号单元，\texttt{>} 右移，\texttt{,} 读次个字符入 1 号单元，\texttt{<} 回到 0 号单元；\texttt{[->+<]} 反复把 0 号单元减 1、1 号单元加 1，直至 0 号单元归零，于是 1 号单元变成两字符 ASCII 值之和（首个字符为 0 时循环体被跳过，1 号单元本就是第二个字符，结果仍是两者之和）；最后 \texttt{>.} 右移并输出。

单元值越界按 256 回绕，故和超过 255 时输出的是“和 $\bmod\ 256$”。

\appendix

\section{显存—系统内存回退管线}
\label{app:vram-fallback}
\subsection{结论}

此处不同大模型给出的不同回答异质化程度高\cite{bao2026reliability}，为了确证显存不足时能否回落到持久化存储层，对该问题进行文档查证，结论如下：在 Windows 的 WDDM 模型下，显存不足时系统可以把 allocation 从 VRAM 驱逐到系统内存；这部分系统内存并不是“第二块显存”，只有被驱逐后处于 non-resident 状态的 allocation 才可能进一步由 Windows 分页到 pagefile。GPU 正在访问的 resident allocation 必须驻留物理内存。因此，pagefile 是 WDDM 的最后一级后备存储，而不是 GPU 可直接寻址的显存扩展。

不同机制的能力边界如表~\ref{tab:vram-fallback} 所示。

\begin{table}[htbp]
  \centering
  \small
  \caption{不同机制下显存回退与磁盘分页能力}
  \label{tab:vram-fallback}
  \begin{tabular}{@{}p{0.22\textwidth}p{0.30\textwidth}p{0.38\textwidth}@{}}
    \toprule
    机制 & 显存不足时能否转到系统内存 & 系统内存部分能否再到分页文件/交换区 \\
    \midrule
    Windows WDDM（NVIDIA GeForce 的 CUDA Sysmem Fallback / 共享 GPU 内存）
      & 能。VidMm 将 allocation 驱逐到 system memory segment。
      & 有条件能。只有 non-resident/evicted allocation 的内容可以被 paged out to disk；resident 数据必须在物理内存中。 \\
    \addlinespace
    Linux NVIDIA 独显
      & 不能。驱动不支持 CUDA Sysmem Fallback。
      & 不适用。 \\
    \addlinespace
    Linux CUDA Unified Memory（full support）
      & 能 oversubscribe，页面在 GPU 与 CPU 之间迁移。
      & 普通 UM 页不能 swap 到磁盘，上限是系统 RAM；file-backed UM 是例外。 \\
    \addlinespace
    Windows/WSL CUDA Unified Memory
      & 不能 oversubscribe（limited UM）。
      & 不适用。 \\
    \addlinespace
    集显/UMA
      & 本来就是使用系统内存。
      & 普通系统页可被 OS 换出，但 GPU 访问期间会 resident。 \\
    \bottomrule
  \end{tabular}
\end{table}

\subsection{Windows/WDDM：可以到 pagefile，但只能间接到达}

Microsoft 的 WDDM segment 模型将 GPU 可访问的内存分为 VRAM（memory segment）、aperture segment 与 system memory segment。GPU 页表既可以指向 VRAM，也可以指向系统内存：

\begin{quote}
\small
``The system memory segment is an implicit segment representing system memory references \ldots\ VidMm implicitly enumerates it and always assigns \texttt{SegmentId==0}.''
\end{quote}

\noindent
来源：\href{https://learn.microsoft.com/en-us/windows-hardware/drivers/display/gpu-segments}{Microsoft Learn: GPU Segments}。

当 VRAM 不足时，VidMm 可以把 allocation 从 VRAM 驱逐到系统内存：

\begin{quote}
\small
``If the GPU resources that are associated with an allocation currently mapped for direct application access are evicted, the content of the allocation is transferred to system memory so that the application can continue to access the content at the same virtual address but different physical medium.''
\end{quote}

\noindent
来源：\href{https://learn.microsoft.com/en-us/windows-hardware/drivers/display/mapping-virtual-addresses-to-a-memory-segment}{Microsoft Learn: Mapping Virtual Addresses to a Memory Segment}。

关键在 residency。Microsoft 的 D3D12 residency 文档明确指出：

\begin{quote}
\small
``GPUs do not yet support page-faulting, so applications must commit data into physical memory while the GPU could access it.''

``Heaps can be evicted from residency, which allows their content to be paged out to disk.''

``\ldots the kernel can shift some heaps on discrete adapters from video memory to system memory, it does so only as an extreme last resort.''
\end{quote}

\noindent
来源：\href{https://learn.microsoft.com/en-us/windows/win32/direct3d12/residency}{Microsoft Learn: Residency (D3D12)}。

因此，实际管线是
\begin{center}
\texttt{VRAM} $\rightarrow$ 系统内存（shared GPU memory）$\rightarrow$（被驱逐后）pagefile/磁盘后备存储，
\end{center}
而不是
\begin{center}
\texttt{VRAM} $\rightarrow$ pagefile（GPU 直接访问）。
\end{center}

pagefile 只是后备存储。GPU 不能直接执行或访问磁盘上的页；要访问时必须先 page in 到物理 RAM，再传到 VRAM。WDDM 所说的“分页”是 allocation eviction/residency，而不是 CPU MMU 那种按需页错误。

NVIDIA 的 CUDA Sysmem Fallback 本质上是 WDDM 之上的一层策略开关。NVIDIA 开发者论坛的 Robert Crovella 解释过：

\begin{quote}
\small
``The WDDM system actually manages GPU memory (plus sysmem as a fallback) via something like a demand-paged virtual memory system \ldots\ paging (i.e.\ movement/relocation) of data is possible, between device memory and sysmem.''
\end{quote}

\noindent
来源：\href{https://forums.developer.nvidia.com/t/cudamalloc-with-sysmem-fallback/347791}{NVIDIA Developer Forums: cudaMalloc with sysmem fallback}。

也就是说，真正执行 eviction 与 paging 的是 Windows 的 VidMm，而不是 NVIDIA 驱动；NVIDIA 只提供“是否允许 sysmem fallback”的策略。

实际影响有两点。第一，如果 pagefile 被禁用或系统 commit 上限不足，到磁盘这一级就没有空间；系统内存耗尽后 allocation 会失败或触发 OOM，不会无限回退。第二，即使能够落盘，性能也会极差：VRAM 到系统内存已经是 Microsoft 所说的 ``extreme last resort''，再到 pagefile 只会造成 thrashing，实际表现通常是卡死、TDR 或帧率崩溃，不能当作显存扩展使用。

\subsection{CUDA Unified Memory：普通 UM 页不会 swap 到磁盘}

如果讨论的是 \texttt{cudaMallocManaged} 的 Unified Memory，结论有所不同。

NVIDIA 官方 Programming Guide 只说明 oversubscription 与页面迁移，并未把“是否 swap 到磁盘”写成规范。NVIDIA 论坛的 Robert Crovella 明确回答过：

\begin{quote}
\small
``AFAIK the UM system pages cannot be swapped out to disk. One indicator of this is that the upper bound for GPU oversubscribed allocations is the system memory size \ldots''

``AFAIK, UM pages are never swapped out to disk. The general UM behavior on linux for Pascal or newer GPUs is one of migration.''
\end{quote}

\noindent
来源：\href{https://forums.developer.nvidia.com/t/unified-memory-is-pageable-it-can-be-swap-out-to-disk/328918}{NVIDIA Developer Forums: Unified Memory is pageable?}。

\Needspace*{6\baselineskip}
因此：\keepwithnext
\begin{itemize}
  \item \textbf{Linux full UM}：允许 oversubscribe，页面在 GPU 与系统内存之间迁移；但普通 UM 页的上限是系统 RAM，不会落到 swap。
  \item \textbf{例外}：file-backed UM（\texttt{mmap} 的文件，配合 HMM/ATS）本来就是文件后备，可以被回收；但这不是“swap”，而是文件页缓存或文件后备。
  \item \textbf{Windows/WSL}：CUDA Programming Guide 明确说明这是 limited UM，不能 oversubscribe：
  \begin{quote}
  \small
  ``GPU Memory Oversubscription: They cannot allocate more managed memory than the physical size of GPU memory.''
  \end{quote}
  来源：\href{https://docs.nvidia.com/cuda/cuda-programming-guide/04-special-topics/unified-memory.html\#unified-memory-on-windows-wsl-and-tegra}{CUDA Programming Guide: Unified Memory on Windows, WSL, and Tegra}。
\end{itemize}

需要说明的是，NVIDIA 官方文档没有把 UM 的 swap 行为写到页级；上述判断来自 NVIDIA 开发者论坛的答复，应视为实现说明，而不是正式 API 契约。

\subsection{Linux NVIDIA 独立显卡：没有 sysmem fallback}

NVIDIA 开发者论坛对“Linux 上如何开启 System Memory Fallback”的回复是：

\begin{quote}
\small
``Sorry, not supported by the nvidia linux driver.''
\end{quote}

\noindent
来源：\href{https://forums.developer.nvidia.com/t/system-memory-fallback-for-linux/296874}{NVIDIA Developer Forums: System Memory Fallback for Linux}。

因此，Linux 下没有 VRAM 到系统内存的自动回退；不要指望 swap 能顶显存；CUDA OOM 就是 OOM，只能通过减小模型、量化、CPU offload 或多卡来解决。

\subsection{集显/UMA：本来就是系统内存}

集显没有独立 VRAM，GPU 通过 WDDM 的 IoMmu 模型直接共享 CPU 页表、访问系统内存：

\begin{quote}
\small
``In the IoMmu model, both the CPU and GPU share a common address space and CPU page tables. Only system memory can be accessed in this case, so IoMmu is suitable for integrated GPUs.''
\end{quote}

\noindent
来源：\href{https://learn.microsoft.com/en-us/windows-hardware/drivers/display/gpu-virtual-memory-in-wddm-2-0}{Microsoft Learn: GPU Virtual Memory in WDDM 2.0}。

这些页本质上就是普通系统页，OS 可以换出；但 GPU 正在访问时会被 resident/pin。它不属于“VRAM 回退”，而是本来就用系统内存。

\subsection{实践含义}

\begin{itemize}
  \item \textbf{Windows + 独显}：pagefile 理论上是 WDDM 的最后一级后备存储，但不要把它当显存用。要运行大于 VRAM 的模型，优先使用 CPU offload、量化、减小 batch，而不是加大 pagefile。
  \item \textbf{Linux + NVIDIA 独显}：没有 sysmem fallback，swap 也不会顶 VRAM；直接按 VRAM 上限做规划。
  \item \textbf{CUDA UM}：Linux full UM 可以 oversubscribe 到系统 RAM，但不会到磁盘；Windows limited UM 不能 oversubscribe。
\end{itemize}

\subsection{参考资料}

\begin{itemize}
  \item \href{https://learn.microsoft.com/en-us/windows-hardware/drivers/display/gpu-segments}{Microsoft Learn: GPU Segments}
  \item \href{https://learn.microsoft.com/en-us/windows-hardware/drivers/display/gpu-virtual-memory-in-wddm-2-0}{Microsoft Learn: GPU Virtual Memory in WDDM 2.0}
  \item \href{https://learn.microsoft.com/en-us/windows-hardware/drivers/display/mapping-virtual-addresses-to-a-memory-segment}{Microsoft Learn: Mapping Virtual Addresses to a Memory Segment}
  \item \href{https://learn.microsoft.com/en-us/windows/win32/direct3d12/residency}{Microsoft Learn: Residency (D3D12)}
  \item \href{https://forums.developer.nvidia.com/t/cudamalloc-with-sysmem-fallback/347791}{NVIDIA Developer Forums: cudaMalloc with sysmem fallback}
  \item \href{https://forums.developer.nvidia.com/t/unified-memory-is-pageable-it-can-be-swap-out-to-disk/328918}{NVIDIA Developer Forums: Unified Memory pageable / swap}
  \item \href{https://forums.developer.nvidia.com/t/system-memory-fallback-for-linux/296874}{NVIDIA Developer Forums: System Memory Fallback for Linux}
  \item \href{https://docs.nvidia.com/cuda/cuda-programming-guide/04-special-topics/unified-memory.html}{CUDA Programming Guide: Unified Memory}
\end{itemize}

\section{素材进 git 不用 LFS 的依据}
\label{app:vendor-lfs}

是否需要 LFS 的判断标准是「文件大小 $\times$ 变更频率」，不是「是否二进制」。git 的 delta 压缩对二进制无效，每次\emph{修改}都全量新增一个 blob，仓库因此膨胀——但内容寻址存储对\emph{不变}的文件恰好只存一份。

git 本身对文件大小没有任何限制，常被引用的 50\,MB 警告与 100\,MB 拒收是 GitHub 的托管政策（GitLab 同为 100\,MB），在自建 git 服务里都是可配置项（\href{https://docs.github.com/en/repositories/working-with-files/managing-large-files/about-large-files-on-github}{GitHub Docs}）。GitHub 自家的仓库体检工具 \texttt{git-sizer}（\href{https://github.com/github/git-sizer}{github/git-sizer}）对超限对象的措辞也只是 ``not wrong per se''——本身不算错，超出规模甜蜜区才有性能代价。真正的反模式只有三类：频繁改写的大文件、可再生的生成物、压缩归档（归档之间无法 delta）。

本仓库的素材正是理想的情形：提交后不可变（答卷材料不会改动）、总量 5.2\,MB（远低于上述警告线与 GitHub 建议的 1\,GB 仓库规模）、更新频率趋近于零，增长模型线性且可预测。

用 LFS 反而更差。它是 git 核心之外的第三方方案，git 维护者无意将其收入核心；指针一旦写进对象库即成单向门，采纳或移除都必须重写历史（\href{https://gregoryszorc.com/blog/2021/05/12/why-you-shouldnt-use-git-lfs}{Szorc 2021}）；对所有协作者多一套安装与出错面，免费档配额仅 1\,GB 存储/月流量（\href{https://docs.github.com/en/repositories/working-with-files/using-large-files/about-git-large-file-storage}{GitHub Docs}）。而 git 核心的 partial clone（\texttt{clone --filter=blob:limit=100k}）已原生覆盖「大文件按需拉取」这一 LFS 的核心收益，实测能把大仓库克隆从 3 分 49 秒缩到 6 秒；2025 年起 git 上游还在合入 large object promisors——官方路线就是原生解决大文件、让 LFS 过时（\href{https://tylercipriani.com/blog/2025/08/15/git-lfs/}{Cipriani 2025}）。AOSP 用普通 git 的 \texttt{prebuilts/} 装着几十个版本的 clang 工具链（\href{https://android.googlesource.com/platform/prebuilts/clang/host/linux-x86}{AOSP prebuilts}），是同一判断的工业级先例。

本仓库 \texttt{.gitignore} 排除了真正不该进 git 的东西——\texttt{build/}（可再生）与 zip 交付包——素材与生成物的边界因此清晰。

\subsection{参考资料}

\begin{itemize}
  \item \href{https://docs.github.com/en/repositories/working-with-files/managing-large-files/about-large-files-on-github}{GitHub Docs: About large files on GitHub}——50\,MB/100\,MB 为托管政策而非 git 限制；仓库规模建议（1\,GB 以下，强烈建议 5\,GB 以下）
  \item \href{https://docs.github.com/en/repositories/working-with-files/using-large-files/about-git-large-file-storage}{GitHub Docs: About Git Large File Storage}——LFS 免费档 1\,GB 存储/月流量配额
  \item \href{https://github.com/github/git-sizer}{github/git-sizer}——GitHub 官方仓库体积体检工具
  \item \href{https://gregoryszorc.com/blog/2021/05/12/why-you-shouldnt-use-git-lfs}{Gregory Szorc: Why You Shouldn't Use Git LFS}——LFS 是核心外权宜方案，进出都是单向门
  \item \href{https://tylercipriani.com/blog/2025/08/15/git-lfs/}{Tyler Cipriani: The future of large files in Git is Git}——partial clone 实测（3 分 49 秒 $\rightarrow$ 6 秒）与 large object promisors 路线
  \item \href{https://android.googlesource.com/platform/prebuilts/clang/host/linux-x86}{AOSP: prebuilts/clang}——普通 git 装几十版工具链的工业先例
\end{itemize}

\bibliography{refs}

\end{document}